\section{The large-prime estimate}
\label{sec:tail-proof}

We now assemble the preceding estimates to prove
Proposition~\ref{prop:large-prime-tail}. Throughout this section $f$ is
fixed, primitive and irreducible, with positive leading coefficient,
$4\le d\le8$, and $k=d-2$. For sufficiently large $X$,
\[
 0<f(n)\ll_f X^d\qquad(X<n\le2X).
\]
If $p^k\mid f(n)$, then $p\ll_f X^{d/k}$. Partition the primes $p>X$
into $O_f(\log X)$ dyadic intervals $[P,2P]$, truncating the final interval
as necessary. For each interval, choose one qualifying prime for each
qualifying input. Put
\[
 \eta=\frac{\log P}{\log X},\qquad b=d-k\eta.
\]
Thus $1\le\eta\le d/k+O_f(1/\log X)$. We seek one common positive
saving throughout these intervals. Choices of charts, ideal classes
and components made below have bounded multiplicity and may be
summed at the end.

\subsection{The range handled by the original surface}

Suppose first that
\begin{equation}
 \eta\left(1-\frac{k\eta}{d}\right)\le\frac{2495}{10000}.
 \label{eq:tail-direct-range}
\end{equation}
The chosen triples
\[
 (x,y,z)=\left(n,\frac{f(n)}{p^k},p\right)
\]
are integer points on the irreducible affine surface
\begin{equation}
 f(x)=yz^k.
 \label{eq:tail-direct-surface}
\end{equation}
Its irreducibility follows from Gauss's lemma, since it is primitive
and linear in $y$ over $\C[x,z]$: the coefficients $z^k$ and $f(x)$
are coprime. The coordinate bounds are $O_f(X)$, $O_f(X^b)$ and
$O_f(X^\eta)$.

Use the finite positive rational upper weights $b',\eta'$ constructed
at the end of Section~\ref{sec:parameters}. They satisfy
\begin{equation}
 \sqrt{b'\eta'/d}<0.4999.
 \label{eq:tail-direct-margin}
\end{equation}
The construction covers $b=0$ and the possible
$O_f(1/\log X)$ excess beyond $d/k$, for all sufficiently large $X$.
Work separately with these finitely many choices.

For weights $(1,b',\eta')$, the space of polynomial functions on
\eqref{eq:tail-direct-surface} has filtered dimension
\begin{equation}
 H(T)=\frac{D}{b'\eta'}\frac{T^2}{2}+O(T),\qquad
 D=\max\{d,b'+k\eta'\}.
 \label{eq:tail-surface-hilbert}
\end{equation}
Indeed multiplication by the defining polynomial increases weight
by exactly $D$, since its weighted highest part is nonzero and the
polynomial ring is a domain. Choose a positive integer $L$ so that
$A=L$, $B=Lb'$, $C=L\eta'$ and $E=LD$ are integers.
Weighted homogenization, with a variable of weight one, gives
the filtered Hilbert series
\[
 \frac{1-t^E}{(1-t)(1-t^A)(1-t^B)(1-t^C)}.
\]
The pole at one has order three and leading coefficient $E/(ABC)$.
Every other root of unity has pole order at most two: if its order
divides all three coordinate weights $A,B,C$, it also divides the
monomial weight $E$, so the numerator cancels a denominator factor.
Extracting the coefficient of $t^{\lfloor LT\rfloor}$ and rescaling
gives \eqref{eq:tail-surface-hilbert}, with the stated $O(T)$ error.
Consequently the surface threshold in
Theorem~\ref{thm:affine-counting} may be taken to be
\[
 \delta_2=\sqrt{b'\eta'/D}<0.4999.
\]

Consider an irreducible curve defined over $\Q$ on this surface and
meeting the selected points. A constant $x$ contributes at most one
selected triple. Otherwise $\C(x)$ is a subfield of its function
field. If the extension has degree at least two, one of the other
coordinate functions lies outside $\C(x)$. Multiplying powers of
$x$ by that function and by one gives two independent families,
so $H(T)\ge2T-O(1)$.

In the remaining case $y$ and $z$ are rational functions of $x$.
They are defined over $\Q$: the curve is defined over $\Q$ and the
functions are uniquely determined by $x$. Their degrees are bounded
in terms of the curve degree, by projection. A nonpolynomial
function is covered by Lemma~\ref{lem:rational-values} and
contributes $O_\eps(X^\eps)$ points. If both are polynomials,
the identity $f(x)=y(x)z(x)^k$ and the absence of repeated roots
of $f$ force $z$ to be constant. At selected points this constant
is a prime $p>X$. Outside the fixed bad primes, $f$ has at most
$d$ roots modulo $p^k$, and $p^k>X$; hence this curve contributes
$O_f(1)$ inputs in $(X,2X]$. These alternatives are uniform for
bounded curve degrees.

An irreducible two-dimensional subvariety of
\eqref{eq:tail-direct-surface} is the surface itself. Its Hilbert bound
and the curve alternatives above, including the bounded constant-$x$
and constant-$z$ cases, therefore verify the hypotheses of
Theorem~\ref{thm:affine-counting} with $\delta_1=1/2$.
Since $\delta_2<\delta_1$, that theorem bounds the number of triples by
\[
 O_{f,\eps}\bigl(X^{\delta_2+1/2+\eps}\bigr).
\]
The finite list of choices in \eqref{eq:tail-direct-margin} permits
one sufficiently small $\eps>0$ and one positive power saving.

\subsection{The number-field range}

For all other $\eta$, use one of the finite parameter intervals
from Proposition~\ref{prop:parameter-certificate}. Denote its
parameters by $t_*,w,m,u,v$ and use the thresholds
$\delta_h(U,V)$ defined there. In particular,
\[
 m=(d-1)w,\qquad 0<m<\min(\eta_-,b_-),\qquad
 u=\frac{\eta_+-m}{d},\quad
 v=\frac{b_+-m}{d},\quad \tau=\frac1{50000}.
\]
The finite interval list covers the complement of
\eqref{eq:tail-direct-range}; its stopping endpoints lie strictly
inside the latter range. This supplies coverage at all boundaries.

For $d=4$, remove once the set of exceptional values in
Proposition~\ref{prop:quintic-discard}. Its total size is
$O_f(X^{1/5})$, independently of the subsequent prime intervals,
boxes or curves. Lemma~\ref{lem:balanced-factorization} and
Lemma~\ref{lem:paired-boxes} place the remaining chosen inputs
in paired arithmetic boxes. The strict inequalities in the
parameter certificate permit the cuts of
Proposition~\ref{prop:archimedean-cuts}.

For each box pair, the resulting affine sets have dimension at
most $d-1$, except for the mixed-cut components of dimension $d$.
All their degrees and their number per pair are bounded before
$X$ varies. Group the pairs as in Proposition~\ref{prop:lattice-boxes}.
There are only boundedly many indices $j\ge0$, and a group has
\begin{equation}
 O_f(X^{m-dj})\text{ pairs},\qquad
 |x'_i|\ll_f X^U,\quad |y'_i|\ll_f X^V,
 \quad U=u+j+\tau,\quad V=v+j+\tau.
 \label{eq:tail-lattice-group}
\end{equation}
The coordinate changes are integral and unimodular. Their
coefficients may vary; the preceding Hilbert statements are
uniform under these changes and fixed degree bounds.

In quartic mode $\mathsf Q$, first delete the infinite-fiber
points of Lemma~\ref{lem:quartic-fibers} from each initial set.
The cost per box is $O_f(X^{2U})$. On the remaining point set,
Proposition~\ref{prop:weighted-hilbert},
Proposition~\ref{prop:quartic-surface} and
Proposition~\ref{prop:curve-alternatives} give, for every
irreducible $\Q$-defined subvariety that can occur in the affine
counting theorem, either its uniform small-count alternative or
\[
 H_A(T)\ge\frac{T^h}{h!\,\delta_h(U,V)^h}-O(T^{h-1}),
 \qquad h=\dim A>0.
\]
Outside mode $\mathsf Q$ the general Hilbert and curve statements
give the same assertion with the corresponding thresholds. A
constant input or constant first coordinate tuple has the already
proved bounded count. The deletion of the quintic values is
precisely what supplies the improved quartic curve threshold.

Set $E_h(U,V)=\max_{h\le s\le d}\delta_s(U,V)$.
Theorem~\ref{thm:affine-counting}, with ambient dimension $2d$
and dimension bound $d$, bounds each remaining set by
$O_{f,\eps}(X^{\sum_hE_h(U,V)+\eps})$.
Its hypotheses have been checked for every bounded degree, so
they apply to the moving subvarieties introduced by that theorem,
not just to the initial cut components. Multiplying by
\eqref{eq:tail-lattice-group}, and applying the shift inequality
in the parameter certificate, gives the exponent
\begin{align*}
 m-dj+\sum_{h=1}^d E_h(u+j+\tau,v+j+\tau)+\eps
 &\le m-dj+\sum_{h=1}^d\bigl(E_h(u,v)+j+\tau\bigr)+\eps\\
 &=m+\sum_{h=1}^dE_h(u,v)+d\tau+\eps.
\end{align*}
By \eqref{eq:parameter-final-box-exponent}, this is less than
$24979/25000+\eps$. A sufficiently small fixed $\eps$ therefore
leaves a uniform positive gap below one.

The quartic infinite-fiber cost, including the number of boxes, has
exponent
\[
 m-dj+2(u+j+\tau)<3122/3125<1
\]
by \eqref{eq:parameter-fiber-exponent}. Thus it also has a uniform
positive saving. The finite number of
parameter intervals, boundedly many lattice groups and bounded
component multiplicities affect only the implied constant.

\subsection{Completion}

The determinant degrees, patch degrees, Hilbert error constants
and counting tolerances are fixed using the strict margins before
$X$ tends to infinity. There are finitely many parameter intervals;
take the minimum of their savings and that of the direct surface
range. Increasing one threshold for $X$ absorbs all bounded
constants. Summing over the $O_f(\log X)$ dyadic prime intervals
then preserves a smaller positive power saving. The single
quartic exceptional set is also smaller than this bound after
decreasing the saving if necessary. This proves
Proposition~\ref{prop:large-prime-tail}.

Finally, Proposition~\ref{prop:sieve-transfer} applies to a
sign-normalized primitive part of each polynomial in
Theorem~\ref{thm:density}. It retains the exact local factors of
the original polynomial, including its content primes, and gives
the asserted positive Euler-product asymptotic.
